% !TeX root = ../papers/nc-p3.tex
\section{種数3の全通常点の四次中心と有限剰余}
\label{sec:rank-fifteen-octic-center}
\begin{theorem}\label{thm:rank-fifteen-octic-center}
階数15の二表示で、$\xi^2+\xi+1=0$ の両水準の全通常点が文字別四次中心次元
$(3,2,0,0)$ を持つ。以下の関数体消去で残る7104幾何点は
標準Koszul左右AS正則である。
\end{theorem}
\begin{proof}
八次式を $F=t^8+14t^4+1$ とする。$i^2=-1$ として
$E=\mathbb Q(i)(t,Y)$、$Y^2=F$ を取り、さらに $d^2=3$ とする。
$j=(d+i)/2$ は $j^4-j^2+1=0$ を満たす。
$\mathbb Q(i,d)=\mathbb Q(j)$ なので $E(d)$ における一つの消去で
全ての定数共役を扱える。$F$ は平方自由で、定数体を拡げても非平方である。
\path{computations/rank_fifteen_octic_field.py} は各係数を $A+dB$、$A,B\in E$ として
持ち、$d$ についてのノルム $A^2-3B^2$ で逆を取る。八次モデルの $K,H,r,w$ を使うと
$(a,b,e,f)=(1+H/j,(K+jH)/r,1-H/j,(K-jH)/r)$ となる。

\path{computations/verify_rank_fifteen_octic_center.py} は長さ5までの道を左に延ばす。
109の $(t,Y,i,d)=(2,46,33,49)$ で独立な行・ピボット列を選び、同じ消去を
標数零の上記の体で行う。各段で残る語の数は既知の経路次元と一致するので、
選んだ行だけで必要な全行の空間を張る。最終射影の係数 $A,B$ の分母と関係式の
分母を記録する。定理\ref{thm:fixed-opposite-rank-drop}で長さ4,5の経路次元は
各通常点で一定である。このため射影が正則な点では、恒等的に関係を殺し
自由語をそのまま返す射影を特殊化して実際の経路基底を得る。

各文字の交換子行列の独立な行を同様に取り、得られる小行列式のノルムを計算する。
具体的にはピボットごとに $E(d)/E$ と $E/\mathbb Q(t)$ のノルムを順に取り、
分子・分母を交叉約分して掛ける。得られる階数での核次元は $3,2,0,0$ である。
一般点での同じ下界は、定次元の経路上で中心の核を特殊化すれば各点でも成り立つ。
射影の極と小行列式の例外因子を合わせ、前節までに求めた非通常因子を除くと、
各表示につき17個の相異なる第一座標の既約因子が残る。双方でその次数の多重集合は
\[
 4,4,8,8,16,16,16,16,48,64,80,80,80,80,96,136,136
\]
である。

各既約因子が単根を持ち、$j$ とその $t$ 上の $Y$ が全て単純分解する素数を取る。
各素体で全四根 $j$ と両根 $Y$ において道、中心および四式の中心商を取り直すと、
全四始点で長さ13まで標準列を得る。単根の持上げおよび共役から、対応する通常点の
全次数の正則性が従う。一つの $\xi$ に $j$ の選択を固定すれば $t,Y$ の表示は
通常域で逆を持つ。従って二表示、二水準、次数和888、二枚の $Y$ を数えて7104点となる。
\end{proof}

\begin{proposition}\label{prop:rank-fifteen-octic-residues}
この四曲線の通常な標数零窓で、109での射影的な $t$ の還元が
$\mathbb P^1(\mathbb F_{109})$ に属し、
\[
 \bar t\notin\{5,22,33,37,53,56,72,76,87,104\}
\]
となるものは標準Koszul左右AS正則である。また $K^2+K-1$ が109で零に
還元される全通常窓も、$t$ の剰余次数によらず同じ正則性を持つ。
\end{proposition}
\begin{proof}
八次モデルの境界の第一座標を109で分解する。$K=\infty$ に対応する
$\bar t=0,1,-1,\infty$ は定理\ref{thm:rank-fifteen-pencil-infinity}で認証済みである。
残る有限の非通常有理剰余は命題に記した十個で、通常な96個は
$t\mapsto-1/t$ で48対になる。\path{computations/verify_rank_fifteen_octic_residues.py}
は各対の一代表で両表示を検算する。$Y$ の二根が素体にあれば両方、
非分解ならFrobenius共役の一方を取る。$j$ は符号を除いて $8,41$ で十分である。
実際 $j\mapsto-j$ は $H,r,w$ を保ち $n$ を反転し、窓の二位置移動を与える。
各代表では全四始点の中心商を長さ13まで確認した。

最後の境界は $K=10,98$ に分かれる。このとき $r=\pm i$、$w=2K+v$、
$v^2=4K^2-1$、$H^2=K^2-1$ である。$z=-1$ だが、使う関係の表示は正則である。
実際その分母に現れる $z,\lambda,z-1,\lambda^2 z-1$ は全て単元であり、
剰余関係は付値環上の通常窓の二次関係を特殊化して得られる。
$H,v$ はいずれも109で非平方の元の平方根で、Frobeniusは両符号を同時に反転する。
$K,r,w$ を固定して $H$ だけを反転すれば $n$ が反転して二位置移動になる。
従って $r$ の二値、それぞれの $j$ の符号代表、両表示について一組の $H,v$ を
検査すれば全選択を尽くす。ここでも隣接階数6、経路次元35,56、中心次元および
長さ13までの商の標準列が成立する。中心が持ち上がるので、通常な各付値体の
窓では全次数の自由性と正則性が従う。
\end{proof}
